\section{Robustness and Discussion}\label{sec:discussion}

The empirical results should be interpreted through robustness rather than rhetoric. LightGBM remains the most credible baseline, whereas Ridge achieves larger gaps only with much poorer fit and the stacked LSTM fails to preserve a convincing average post-policy signal. The rolling-origin check also keeps LightGBM preferred, although $w=6$ and $w=12$ are close enough that the exact short-range window should not be overinterpreted.

The added diagnostics strengthen but do not settle the case. The preferred model beats the naive baselines, retains positive bootstrap intervals, and its 60-day January 2022 boundary is stronger than the nearby October and November 2021 placebos. Yet the placebo cuts are not null. The placebo contrast is supportive but not dispositive, and the forecasting design attenuates temporal drift rather than eliminating it.

The traffic and evaluation sensitivities mainly clarify interpretation. Removing meteorology weakens the signal more than removing district traffic intensity, so the result should be read as a deviation after adjustment for weather and available urban covariates, not as identification of a traffic mechanism. The small MAE advantage of the no-traffic ablation is therefore diagnostic rather than dispositive: it shows limited incremental predictive contribution from the available district proxy, not that mobility information is substantively irrelevant. Likewise, a homogeneous direct one-step evaluation keeps LightGBM best by MAE but attenuates the gap once observed post-policy lags enter the window. That is exactly what should happen: the empirical question changes once treated lags replace the untreated autoregressive path. Recursive rollout therefore remains necessary for the no-change estimand.

The district narrative and the external-validity claim must therefore stay local. The strongest positive cell is PM$_{10}$ in Villa de Vallecas, whereas NO remains slightly negative in the three districts. This is evidence of local heterogeneity in the selected districts and retained stations, not proof of a city-wide causal effect.
